\documentclass[12pt, twoside, a4paper]{report}

\usepackage[bottom=1 cm, top=1 cm, left=1.8 cm, right=1.8 cm, includehead, includefoot]{geometry}
\usepackage{titling}
\usepackage{fontspec}
%\setmainfont{Liberation Sans}
\setmainfont{Linux Biolinum}
%\setmonofont{Fira Code}
\setmonofont{Nimbus Mono PS}
\usepackage{array}

\usepackage{tikz}
\usetikzlibrary{arrows.meta}
\usepackage{xcolor}

\usepackage{listings}
\lstset{
	basicstyle=\ttfamily,
	numbers=left,
	numberstyle=\tiny,
	frame=tb,
	tabsize=4,
	columns=fixed,
	showstringspaces=false,
	showtabs=false,
	breaklines=true,
	%keepspaces,
	commentstyle=\color{gray},
	keywordstyle=\color{orange}
	}


\usepackage[linesnumbered]{algorithm2e}
%\SetArgSty{textnormal}
\SetKwProg{Fn}{Function}{}{end}

\usepackage{comment}
\usepackage{amsmath, amssymb}
\usepackage{graphicx}
\usepackage{svg}
\usepackage{geometry}
\usepackage{bbold}
\usepackage{parskip}
\usepackage{fancyhdr}
\usepackage{siunitx} % probably not worth it actually
\DeclareSIUnit\bp{bp} % bp defined

% configure biblatex, more styles
%\usepackage{biblatex} %Imports biblatex package
%\addbibresource{bibliography.bib} %Import the bibliography file

\usepackage{xcolor}
\definecolor{lb}{RGB}{44, 139, 183}
\newcommand\cl[2]{\textcolor{#1}{#2}}
\newcommand\bcl[2]{\textcolor{#1}{{\fontseries{b}\selectfont #2}}}

\usepackage[colorlinks=true, breaklinks=true]{hyperref}
\hypersetup{colorlinks, breaklinks, urlcolor=[rgb]{0,0.4,0.7}, citecolor=black, linkcolor=black!75}

\newtheorem{theorem}{Theorem}

\newcommand{\HRule}{\rule{.8\linewidth}{0.5mm}}

\newcommand{\beginsupplement}{%
        \setcounter{table}{0}
        \renewcommand{\thetable}{A\arabic{table}}%
        \setcounter{figure}{0}
        \renewcommand{\thefigure}{A\arabic{figure}}%
     }

%\title{$k$-mer-based Reconstruction of Haplotypes using a Pangenomic Reference}
\title{Representing Phased Haplotypes as a Pangenome}
\author{Leon Rauschning\thanks{\href{mailto:leon.rauschning@tum.de}{leon.rauschning@tum.de}}}

\begin{document}

\begin{titlepage}
\centering
% Oberer Teil der Titelseite:
\includegraphics[width=0.3\textwidth]{img/logo2.jpg}
\hfill
\includegraphics[width=0.4\textwidth]{img/logo1.jpg}  
\\[5cm]

\Large
Bioinformatics Programme\\[0.5cm]
Technical University of Munich\\[0.5cm]
Ludwig-Maximilians-Universität München\\[2cm]
Master's Thesis in Bioinformatics\\[1.5cm]


\HRule \\[0.6cm]
{ \LARGE \bfseries \thetitle}\\[0.4cm]

\HRule \\[2.5cm]

Leon Rauschning\\[2.5cm]

\vfill
\end{titlepage}

\begin{titlepage}
\centering
\includegraphics[width=0.3\textwidth]{img/logo2.jpg}
\hfill
\includegraphics[width=0.4\textwidth]{img/logo1.jpg}  
\\[3cm]  
\large
Bioinformatics Programme\\[0.5cm]
Technical University of Munich\\[0.5cm]
Ludwig-Maximilians-Universität München\\[2cm]
Master's Thesis in Bioinformatics\\[2cm]
{\Large\bfseries
\thetitle\\[2cm]
%Representing phased haplotypes as a pangenome
%$k$-mer-basierte Haplotyprekonstruktion anhand einer pangenomischen Referenz\\[3cm]
Repräsentation haplotypaufgelöster Genome als Pangenom\\[3cm]
}

\large
\begin{tabular}{ll}
 Author:& \theauthor\\
 Supervisor: &  Dr.\ Craig Dent\\
 Advisor:        &  Prof.\ Korbinian Schneeberger, Faculty of Biology, LMU\\
 & Prof.\ Mathias Wilhelm, School of Life Sciences, TUM\\
 Submitted:     &  15/04/2026
\end{tabular}
\end{titlepage}

\begin{titlepage}
\centering
\hspace{0pt}
\vfill
\begin{minipage}{.6\textwidth}
\centering
I confirm that this Master's thesis is my own work and that I have documented all sources and materials used.
\\[2cm]
\rule{7cm}{1pt}\\
15th of April 2026, Munich
\end{minipage}
\vfill
\hspace{0pt}
\end{titlepage}


%\thispagestyle{empty}
\chapter*{Abstract}
% abstract classes cannot be instantiated

PS: citations in abstract yay or nay?

%As it is becoming 
Assembling a complete eukaryotic genome from telomere to telomere has gone from demanding substantial efforts from a well-resourced consortium to routine work of a single lab processing many samples since the T2T project's release of the CHM13 genome in 2022. % cite T2T?
% frame phasing problem + interpretation as two remaining challenges
This leaves one remaining challenge in reconstructing the complete genome of an organism -- separating the different copies of the genome, also known as phased assembly.
In human, this is now possible using the same sequencing reads used to assemble the genome, however this has so far not been demonstrated in more complex genomes.
%It is as of yet unclear if this is impossible in principle, 
\\ %, polyploid genomes.
But once a genome is assembled, further challenges await: genomes need to be compared, annotated or otherwise analyzed to derive biological knowledge from them.
Any analysis requires a suitable representation.
Representing modern datasets of many high-quality genomes is an open problem; arguably, representing the genomes of a large dataset is now harder than assembling them.
%and arguably we are further from a solution to it than genome assembley
%While a single genome can be adequately represented in FASTA format, representing multiple genomes is a challenge that is far less solved 

In this thesis, I present contributions to both of these open challenges in genomics.
First, a model of read-based phasing is introduced, and used to derive requirements for a reads dataset to contain sufficient information to phase a genome.
This proves that read-based phasing can be feasible even in polyploids, but current sequencing technology is still unable to fully phase the potato genome as an example of a particularly tough tetraploid.

Second, a graph genome representation termed the Multisyntenic Pangenome Graph is presented.
%To adequately represent large assembly datasets, I propose the Multisyntenic Pangenome Graph (MPG).
It is based on a previously developed extension of the concept of synteny to multiple genomes, and uses this to derive a representation possessing a linear backbone while avoiding cycles and nonlocal edges.
%underlying the SyRI algorithm, and accounts for global genome structure when collapsing sequences into nodes.
%This provides guarantees that other proposed representations do not, in particular it avoids cycles and nonlocal edges, and provides a linear backbone of shared structure across all samples.\\
The representation is shown to have a natural index preserving some of the favourable properties of linear genome indices, demonstrating the utility of a conceptually motivated representation with strong guarantees.
%This allows efficient description and navigation along the graph.

It will no doubt be exciting to watch the future of genomics unfold as genome assembly is increasingly becoming a solved problem.
I anticipate conceptual and linguistic clarity, as well as the production of not just more data but also novel understanding to be crucial, and hope that this thesis can serve as a small step in this direction for the two areas outlined above. %will be crucialfor theory to catch up to the data.

\textit{too philosopical?}


\pagestyle{fancy}
\fancyhead{}
\fancyhead[C]{\thetitle}
\fancyhead[RE,LO]{Leon Rauschning}
\fancyfoot{} % clear all footer fields
\fancyfoot[LE,RO]{\thepage}


%\chapter*{\thetitle}
\chapter*{Introduction}


% how to approach?
% Idea 1:
	% go heavy on pangenomics
	% namecheck major initiatives
	% move away from linear reference
% not sure; pitches very abstract, no clear motivation
% Idea 2:
	% come from potato breeding perspective
	% deleterious variants => need phased haplotypes
	% shift some pangenome stuff into main/methods?
	% go more generally into phased genome assembly etc in intro
% Craig:
% focus on difficulty of de novo assembly
% => pangenie shows this works in diploids, try in polyploids


% why do we need to represent genomes?
% how do we represent genomes?
% => go explicit about language limitation
% => maybe trace history => dobzhansky, cytogenetics

Recent progress in genome sequencing and assembly technologies have made full genome assemblies more and more accessible, to the point where it is now feasible to obtain assemblies for entire sample populations in many species.
Translating this data into knowledge has, however, been a significant challenge so far.
Genome assemblies are a large and heterogeneous kind of data.
Making the information contained within them accessible is far from trivial, and requires computational approaches scaling to efficiently utilize large computational resources.

The matter of how to represent today's genome assembly datasets is absolutely crucial to any attempt at analyzing them, both in allowing efficient algorithms to be constructed and in limiting, enabling or focusing the questions that we ask a dataset.
While for previous resequencing datasets the VCF format has emerged as a de facto standard -- despite significant limitations especially regarding storage efficiency --, there is still considerable disagreement about how to represent modern datasets of genome assemblies. 
Below, a brief outline of how genomes have historically been represented -- both on analog or digital media and in theoretical discussions -- is given, to provide some perspective on the contemporary discussion.

The first recognition of the genome as a coherent concept dates to Gregor Mendels work in peas.
He represented the genome of his peas as combinations of upper/lowercase letters indicating different alleles.
Given the lack of any genomic technologies, the genetic identity of a given pea could only be determined by its phenotype or the phenotypes of its offspring in a controlled cross.
Mendel's genomes also lacked any internal structure; gametes were viewed as produced by of independently assorted components of the parental genotypes.
The genome as such was a passive container of alleles, and not conceptualized beyond this.\\
After the rediscovery of Mendel's work in the early 20th century, the discovery of chromosomes as the bearer of heredity information and further crossing experiments performed among others by Thomas Hunt Morgan's laboratory at Columbia University in \textit{Drosophila melanogaster} revealed the genome to have a linear structure.
Through careful crossing, chromosomes were mapped, again determining the genotype solely from the phenotype of the organism and its offspring.
These genetic maps revealed chromosomes to have a linear structure, and to be shuffled by a process of recombination by discrete crossover events, happening with a determinable likelyhood between genes segregating for certain phenotypes.\\
The genome had evolved from a few letters to a set of linear maps, whose cellular component (if not molecular identity) was known.
Like Mendel, however, genes along the genome were thought of and named for their observable, segregating phenotypes, and the absence of any genomic technologies made genotyping a slow and labourious process.\\
Later, microscopy of chromosomes revealed a banded structure, which work in maize by Barbara McClintock and in \textit{Drosophila spp.} by Theodosius Dobzhansky showed to sometimes vary across individuals, initially without giving rise to observable macroscopic phenotypes.
The discovery of this so-called »cryptic polymorphisms« was the first time a genotype had been determined by other means than phenotypic observation.
Cytogenetics can thus be termed the first genomic technology; in fact, through careful observation of the band structure, these polymorphisms could be identified as inverted or copy-pasted segments along the chromosomes, hinting at the mechanism of their formation -- later giving rise to McClintock's discovery of transposable elements.\\
Around the same time, the synthesis of Mendel's rediscovered works and contemporary evolutionary biology, as well as contacts with ecology led to the emergence of an early theory of population genetics, a first way to systematically think about the behaviour of genomic variation in a population.
Applying this new population genetic theory to the observed allele frequency of inversions in captive \textit{D. pseudoobscura} populations, these »cryptic polymorphisms« were revealed to have an effect on observed fitness, in particular in heterozygotes, a phenomenon that was then termed heterosis.
% expand on this?
Increasingly, the study of the genome became emancipated from the reliance of clearly segregating phenotypes for every genetic polymorphism, and the discovery of rearrangements revealed that a genetic or cytogenetic map could vary between different individuals of a population.
Genotyping technology allowed population genetic studies for the first time, though these studied a limited amount of polymorphisms at a time.

The discovery first of DNA as the molecule carrying genetic information and later of its 3D structure in the 1940s and 1950s by Avery and Franklin, Crick and Watson provided for the first time a vision to understand the genome on a molecular level.
More molecular studies followed like the discovery of transcription, replication and the decoding of the genetic code, providing some understanding on a basepair level.
This opened the door to a radically different way of conceptualizing the genome -- not as an ordered map of segregating phenotypes or bands under a microscope, but as a string of bases along a phosphate backbone.
As DNA became accessible by cloning, PCR and eventually early sequencing methods, databases of DNA sequences started to emerge.
These would contain small fragments of the genome, often genes of particular interest, promotor or plasmid sequences.
The sequences would be stored as letters, but unlike Mendel's early use of letters to indicate alleles, the letters now indicated the bases \texttt{A}denine, \texttt{C}ytosine, \texttt{G}uanine and \texttt{T}hymine and had a well-defined ordering.
Early databases were still printed on paper and distributed by mail as books, but soon started shifting to early computers available at universities.
Databases used a variety of file formats, but for nucleic acid databases the format supported by the fast-All programme, colloquially known as FASTA format, soon took the lead.
This format was analogous to the printed book format, containing the sequence in an ACGT/ACGU-encoding with nonsignificant linebreaks every 80 characters to fit on early computer screens, but added a header line for each sequence containing the sequence name and additional metadata, if present. 
It is still in common use today for a variety of purposes.
In particular, FASTA has become the de facto standard for representing reference genome assemblies after being chosen by the Human Genome Project and other early genome sequencing collaborations to publish its genomes.


% shift over to VCF format, 1000 GP

% short comment on embeddings/neural nets in general

\newpage
\chapter*{Main}



\section*{Phaseability}

% define min. read len as fn of heterozygosity
% => handle polyploids: multiply heterozygous vars?
% => account for pop structure/IBD afterwards
% consrtuct single probability output, then solve for read len or heterozygosity

In this section, we want to obtain minimal input requirements for reconstructing phased potato genomes, and look at how different properties of our sequencing technology and genome characteristics affect these requirements.
To approach this problem, some clarity on what phasing is and is not is required, in particular how it relates to genome assembly.
In the cells of a given biological organism, there are often multiple copies of the genome – depending on species and tissue, 10 or more are not unheard of.
Because of standing variation or spontaneous mutational processes, these so-called haplotypes may differ between another in minor or major ways.
Variation on the same haplotype is inherited together, unless recombination breaks it up; haplotypes are thus also referred to as the units of heredity.
We term a position at which a haplotype is different from all others a haplo-unique site; they may correspond to single small variants, (parts of) larger variants, or clusters of multiple nearby variants.
A phasing of a genome refers to a map taking two different haplo-unique sites, and telling us whether they are found on the same haplotype or not.
How do we construct such a phasing?
We will mainly concern ourselves with the minimal requirements for constructing a phasing from a set of genomic sequencing reads, but it should be noted that HiC sequencing, a mapping population or parental sequencing may also serve to accomplish this goal, though less economically. % justify why only deriving minimal estimates/can assemblers reach this?
When encountering a read, we know all of its sequence to be derived from a single physical piece of DNA corresponding to a fragment of a single haplotype\footnote{With the possible exception of PCR artefacts.}.
If two haplo-unique sequences $u_0, u_1$ are found on the same read $(u_0, u_1)$, we can assign them to the same haplotype; this construction can be transitively extended, if we have two reads $(u_0, u_1), (u_1, u_2)$, we can use the shared haplo-unique sequence $u_1$ to assign $u_0, u_2$ to the same haplotype.
Note that, contrary to genome assembly, repetitive sequences are not inherently problematic; so long as the sequence contains sufficient haplo-unique sequences, it can be easy to phase but difficult to assemble; vice versa, a perfectly nonrepetitive sequence may be difficult or impossible to phase but easy to assemble if it contains large stretches shared by all haplotypes.

We can thus derive a simple minimal requirement for reconstructing the phasing of a genome, given in theorem \ref{theo-phase-simple}

\begin{theorem}
	Given a set of reads $\mathcal{R}$ and a set of haplotypes $\mathcal{H}$, each modelled as consisting of an ordered set of sequences unique to it $\mathcal{H}_i := (u^i_0, u^i_1, \dots u^i_n)$, it is possible to construct a phasing $\mathcal{P}$ iff for every adjacent pair of haplo-unique sequences $u_j, u_{j+1}$ there is a read covering both: $\forall (u_j, u_{j+1}) \exists x \in \mathcal{R}: (u_j, u_{j+1}) \in x$.

\label{theo-phase-simple}
\end{theorem}

%TODO proof
A proof is simple to construct: if two haplo-unique sites $u_i, u_j$ with $i < j$ are on the same haplotype, a chain of neighbouring haplo-unique sequences $u_i, u_{i+1}, \dots, u_{j-1}, u_{j}$ can be found and thus $\mathcal{P}(u_i, u_j)$ evaluates to true.
If we know all neighbouring haplo-unique sequences to be covered, the absence of such a chain implies that $u_i$ and $u_j$ do not share a haplotype, and thus $\matcal{P}(u_i, u_j)$ can safely evaluate to false. 
Conversely, if there exists a pair of haplounique sequences $u_i, u_{i+1}$ on the same haplotype that are not covered by any read, we will not be able to resolve them unambiguously to recover their phasing.

%TODO statmodelling
Using some statistical modelling, we can turn this simple minimal requirement into an estimate of the sequencing coverage required to confidently resolve the phasing of a genome.
To ensure full phasing, the requirement of theorem \ref{theo-phase-simple} has to be met everywhere along the genome.
In particular, this means that
\begin{enumerate}
	\item every haplotype has to be sampled
	\item every position along the genome has to be sampled, and
	\item the read covering every position on every haplotype has to actually cover both adjacent haplo-unique sequences.
\end{enumerate}
We can model each of these requirements in turn, allowing us to derive a sequencing coverage estimate.

To model the first requirement, if we consider a given position along the genome, every read mapping there will be drawn from all possible haplotypes with equal probability.
Sampling is with replacement, if we assume our sequencing library to have sufficient size.
Asking how many samples need to be drawn to cover every haplotype with confidence is equivalent to the Coupon Collector Problem, for which multiple approximations have been previously derived. % find original citation, derive

The second requirement hearkens back to the early days of estimating sequencing requirements for whole-genome shotgun sequences -- except instead of ensuring every position is covered at least once, we need every position to be covered in sufficient depth to make sure the first requirement is met.
If we model observed coverage along the genome given a specified mean coverage using a Poisson distribution, we can derive a requirement for the mean coverage by solving for the quantile of the Poisson distribution.
This corresponds exactly to the inverse of the upper $\Gamma$ function, a numeric approximation of which is available in the \texttt{zipfR} package.

Finally, not all reads covering a certain position may be informative, if they happen to cover just one of the haplo-unique sequences.
We can compute the fraction of informative reads spanning both haplo-unique sequences from the fraction of valid starting positions, given that every read is uniquely specified by its starting position.
A starting position is valid if a read of length $\ell$ starting from it spans both haplo-unique sequences; all valid starting positions therefore have to be before the first haplo-unique sequences, and there will be $\ell-d$ many, with $d$ giving the distance between the two haplo-unique sequences.
As every read has a single starting position, a mean sequencing coverage of $X$ coverage corresponds to $X/\ell$ reads starting at every position along the genome in expectation.
We thus expect $(\ell-d)\frac{X}{\ell}=\frac{\ell-d}{\ell}X$ reads to fulfil this criterion.
As we assume reads to be Poisson distributed, we can simply multiply our coverage estimate derived in step 2 with this factor.

% incorporate after notion of msyd graph construction

\section*{Proof of Non-Circularity}
\newcommand{\aln}{\mathcal{A}}

We consider two multisyntenic regions $M$ and $M'$ that are both upstream and downstream of each other in genomes $q$ and $q'$, forming a circle in the graph representation.
Because downstream/upstreamness is transitive, this case covers all possible circles in the graph.
Formally, we can express this as $\exists q, q': M[q] < M'[q] \wedge M'[q'] < M[q']$.
Consider the alignment $\aln_{q \to q'}$.
If this hasn't explicitly been constructed, we can induce it from two alignments to a common reference $r$: $\forall i \aln_{q \to q'}:= \aln^{-1}_{q' \to r}[\aln_{q \to r}[i]]$.
Such an $r$ must exist, and be either the reference used during the realignment leading to $M$ or $M'$.
We can quickly see the two parts of the statement leading to a contradiction when we consider that multisyntenic regions are defined to be aligning on all organisms they cover:% at $\aln_{q \to q'}$ , we can see a contradiction emerge quickly to our :

\begin{eqnarray*}
M[q] &< M'[q] &\\
\aln_{q\to q'}[M[q]] &< \aln_{q \to q'}[M'[q]] &\mathrm{Alignments are monotonic}\\
M[q'] &< M'[q'] &\mathrm{but} M'[q'] < M'[q]\mathrm{!}
\end{eqnarray*}

%There are a few assumptions underlying this proof, some of which may be violated in practice.
Note that this argument assumes that our realignment finds monotonic and correct alignments.
\texttt{minimap2}'s fast approximative minimizer-based algorithm is used for computing alignments, which may in practice sometimes violate these conditions.
An additional source of alignment ambiguity is the stitching together of leftover non-multisyntenic sequence during the realignment step.\\
This proof does thus not guarantee that an msyd run may not produce output inducing a circular graph representation in certain edgecases, but instead serve to justify its pruning as making an approximative component of the algorithm match the mathematical guarantees of an exact algorithm.


%\section*{Methods}
%\input{methods}
%\newpage
%\section*{Results} % maybe merge into common main block?
%\input{results}
%\newpage
\chapter*{Discussion}


% pick up on representation from intro?

%% further work, software improvement
% validate covcalc in practice
% ways to improve msyd:
% move realignment to exact library

Recently, Samija Amar in the lab managed to design a protocol to reliably obtain ultralong reads > \qty{200}{\kilo\bp} from potato.
According to the sequencing minima derived in the first section of the main, this should make it feasible to generate read datasets with sufficient information to, in theory, phase all but the largest homozygous regions of the potato genome.%generate be sufficient to 
Testing whether current algorithms can leverage this information for single-platform phased assembly of potato genomes will be interesting, either enabling the high-throughput generation of potato genomes, or providing an opportunity for the development of such tools\footnote{Of note, none of the leading assembly tools were designed from the ground up with ultralong reads in mind; there may be significant potential still left to obtain from these data.}.

A clear next step is to investigate and visualize

Regarding msyd, during dataset construction the realignment functionality of msyd proved to not unreliable, particularly in the potato dataset.
This causes the multisyntenic pangenome graph to be undercollapsed in non-reference haplotypes.
Testing problematic cases revealed that minimap2 may not report any alignment in short sequences, likely as a result of not finding enough minimizers.
Porting the realignment step to use a Dynamic Programming-based aligner, for example Parasail, may make it more reliable\cite{daily16}.
For short sequences, this should not significantly change runtime; for longer sequences, perhaps a fallback to minimap2 or another heuristic approach may still be desirable.

%% Representation stuff
% comparison to other tools from msygraph + panco
% panco
% MSA consensus stuff

Comparing to pangenome graphs based on local-alignments and variation graphs, MPG differ in a number of key ways.
First, small variants not breaking synteny are not represented in the graph structure.
This is similar to a common filtering step used for variation graphs \cite{matthews24, hickey24}.
As the graph can be annotated with exact positions on each sample and aligments are available in the PSF file, these can be recovered in a fairly straightforward way.
Second, a MPG will never introduce a circle by collapsing repetitive sequences, and will also never collapse nonlocal nodes or produce nodes with many nonlocal connections.
Similarly, no connection can be made between nodes on different chromosomes, ensuring that chromosomes can always cleanly be separated out as independent subgraphs.
In fact, the emitted GFA file contains comments indicating the chromosome subgraphs.

At the same time, as syntenic alignments are found using minimap2's approximate and SyRI may also fall back to a heuristic approach in complex situations, the multisyntenic pangenome graph constructed by msyd may in practice be less than fully collapsed, particularly in highly diverged or repetitive regions.
Approaches to further collapse the graph or filter out rare or small nodes may be useful in practice, in particular with regard to private nodes.

Which graph representation is ideal will of course depend on the nature of the dataset and scientific question.
The MPG differs from other current pangenome representations.
It has a more regular structure, and aims to capture large-scale structure over local events.
In the next section, I will outline a procedure using the regular structure of the MPG to define a pangenomic coordinate system termed PanCo, allowing genomic coordinates to be compared across different samples in a principled and consistent way.%\footnote{Note that a similar procedure is in fact used to ensure consistent sorting across samples during MPG construction.}.
This not only serves a direct purpose, but also demonstrates the advantages of a representation possessing strong structural guarantees, and illustrates how these can be used to derive efficient algorithms appropriate for modern genomics datasets.



% discussion stuff?
The fact that this powerful and interpretable coordinate system »falls out« of the graph representation at practically no additional computational cost reflects the advantages of a pangenome representation with strong theoretical guarantees.
%PanCos are not fully stable towards adding additional samples to the pangenome; while the coresyn part of the PanCo will not change, assuming the added sample does not break coresynteny, the merasyntenic part of the PanCo will likely require recomputing.

Compare the ease at which a MPG can be indexed with PanCos to the difficulty of indexing nodes in a variation graph: the odgi sort command implements multiple cycle-breaking algorithms in order to then traverse the entire graph and assign a topological ordering, alongside a force-directed-layout based stochastic embedding approach\cite{guarracino22}.
Meanwhile, I am not aware of any indexing, including via external tools, of minigraph's representations beside arbitrarily assigned numeric node IDs.

% leave comparison/cross species for discussion



%% Future future stuff
% cross-species comparison => panco
% structural site frequency spectrum?



\section*{Acknowledgements}

I would like to thank Craig Dent for his supervision and always helpful feedback.
Further I thank Aleix Canalda-Baltrons, Xiao Dong, Manish Goel, Matthew Hahn, Matthew Parker and Jens Zentgraf for their discussions.
Last, I would like to thank Jumpie for his refreshing distractions while writing.



\pagestyle{plain}
{\small
%\twocolumn
\onecolumn

%biblatex, not bibtex (?)
%\bibliographystyle{nejm}
%\bibliographystyle{nature}

%\bibliographystyle{unsrt}
\bibliographystyle{naturemag}
\bibliography{bibliography.bib}
%\printbibliography
}

{
\onecolumn
\small
\beginsupplement
\newpage
\chapter*{Appendix}
\input{Appendix}
}

\end{document}
